Assume \(\mathcal C(\mathcal P)\ne\varnothing\).  The set \(\Theta(\mathcal P)\) is a nonempty compact semialgebraic set of finite dimension, and canonical realization on the declared observed alphabets gives, for every finite collection \(\mathcal F\) of target-domain counterfactual cells,
\[
 \bigl\{(P_{M^\star}(C):C\in\mathcal F):
       \mathbf M\in\mathcal C(\mathcal P)\bigr\}
 =
 \bigl\{(\Phi_C^\star(\theta^\star):C\in\mathcal F):
       \theta\in\Theta(\mathcal P)\bigr\}.
\]
Consequently, the extrema \(L_T(\mathcal P)\) and \(U_T(\mathcal P)\) are attained and are the sharp bounds for every target cell \(T\).  The bounds are strictly tighter than the logical interval \([0,1]\) exactly when \(L_T(\mathcal P)>0\) or \(U_T(\mathcal P)<1\), an exactly checkable property of the displayed program.  The cell is identified if and only if \(L_T(\mathcal P)=U_T(\mathcal P)\); if \(L_T(\mathcal P)<U_T(\mathcal P)\), minimizing and maximizing canonical parameters yield two SCM families on the original observed alphabets that agree on every cell of every complete supplied law and disagree on \(T\).

Let
\[
 e_-=\min_{\theta\in\Theta(\mathcal P)}\Phi_E^\star(\theta^\star),
 \qquad
 e_+=\max_{\theta\in\Theta(\mathcal P)}\Phi_E^\star(\theta^\star).
\]
If \(e_+ =0\), the evidence has probability zero throughout the compatibility fiber; if \(e_-=0<e_+\), its positivity is not uniform.  In both cases a conditional functional on the entire fiber is rejected before division.  If \(e_->0\), the sharp range of the conditional arm value is
\[
 \left[
 \min_{\theta\in\Theta(\mathcal P)}
 \frac{\Phi_{B_\rho\cap E}^\star(\theta^\star)}
      {\Phi_E^\star(\theta^\star)},
 \quad
 \max_{\theta\in\Theta(\mathcal P)}
 \frac{\Phi_{B_\rho\cap E}^\star(\theta^\star)}
      {\Phi_E^\star(\theta^\star)}
 \right],
\]
with both endpoints attained.  Under the same condition \(e_->0\), uniform strict ordering satisfies
\[
 \begin{array}{rcl}
 Q_\sigma^\star>Q_\rho^\star\text{ in every compatible SCM}
 &\Longleftrightarrow&
 \displaystyle\min_{\theta\in\Theta(\mathcal P)}
 \mathcal G_{\rho,\sigma}(\theta)>0,\\[3pt]
 Q_\rho^\star>Q_\sigma^\star\text{ in every compatible SCM}
 &\Longleftrightarrow&
 \displaystyle\max_{\theta\in\Theta(\mathcal P)}
 \mathcal G_{\rho,\sigma}(\theta)<0,
 \end{array}
\]
and two original-alphabet compatible families with opposite strict rankings exist if and only if
\[
 \min_{\theta\in\Theta(\mathcal P)}\mathcal G_{\rho,\sigma}(\theta)<0
 <\max_{\theta\in\Theta(\mathcal P)}\mathcal G_{\rho,\sigma}(\theta).
\]
In particular, a target structural-causal bandit may safely prune \(\rho\) in favor of \(\sigma\) exactly when the first strict minimum condition holds.

For every district \(D\), the global projection satisfies
\[
 \{q_D^\star(\cdot;\theta^\star):\theta\in\Theta(\mathcal P)\}
 \subseteq \mathcal K_D(m_D).
\]
A local point, or a pair of local points, lifts to globally compatible SCM families exactly when it belongs, or they both belong, to the left-hand projection.  Equality with \(\mathcal K_D(m_D)\) follows under the hypotheses of the one-district fiber-projection theorem, but is not implied in a multi-district problem.  Hence a local \(\mathsf{DCTFTR}\) face-failure trace alone is not a global paired-model certificate; the exact multi-district certificate is instead furnished by distinct optimizers of the displayed global polynomial programs.